
\chapter{Situación inicial y resultados obtenidos}

\section{Situación inicial}

\campo{
Describir el estado del proyecto antes de implementar las
prácticas de DevOps y DevSecOps.
}

\section{Resultados después de la implementación}

\campo{
Describir los resultados obtenidos luego de incorporar
automatización, análisis de seguridad y centralización
de vulnerabilidades.
}

\section{Comparación entre la situación inicial y final}

\begin{table}[H]
\centering
\caption{Comparación del estado inicial y final del proyecto.}
\label{tab:comparacion}

\begin{tabularx}{\textwidth}{l X X X}
\toprule
\textbf{Aspecto} &
\textbf{Situación inicial} &
\textbf{Situación final} &
\textbf{Mejora} \\
\midrule
Integración & \campo{Estado} & \campo{Estado} & \campo{Resultado} \\
Seguridad & \campo{Estado} & \campo{Estado} & \campo{Resultado} \\
Dependencias & \campo{Estado} & \campo{Estado} & \campo{Resultado} \\
Secretos & \campo{Estado} & \campo{Estado} & \campo{Resultado} \\
Vulnerabilidades & \campo{Estado} & \campo{Estado} & \campo{Resultado} \\
Trazabilidad & \campo{Estado} & \campo{Estado} & \campo{Resultado} \\
\bottomrule
\end{tabularx}
\end{table}

\section{Indicadores y métricas}

\campo{
Presentar métricas reales de los análisis realizados.
Cuando no exista una medición inicial, indicar expresamente
esta limitación.
}
